\chapter{Praktikum 1: Akar Persamaan Nonlinear}

\CP{Mahasiswa mampu menerapkan Biseksi dan Regula Falsi sebagai metode pengurungan, memahami Titik Tetap, Newton-Raphson, dan Secant, serta mengevaluasi konvergensi menggunakan residual dan galat iterasi.}

\Target{Kode contoh kelas untuk Biseksi dan Regula Falsi, tabel iterasi, dan grafik konvergensi. Latihan utama dikerjakan dengan Biseksi, Newton-Raphson, dan Secant.}

\section{Masalah pencarian akar}

Tujuan pencarian akar adalah mencari $x$ sehingga

\begin{equation}
f(x)=0.
\end{equation}

Sebelum menghitung, gambar grafik untuk melihat posisi perubahan tanda dan memilih interval atau tebakan awal yang masuk akal.

Urutan metode dalam praktikum ini adalah:

\begin{enumerate}
\item Biseksi,
\item Regula Falsi,
\item Titik Tetap,
\item Newton-Raphson,
\item Secant.
\end{enumerate}

\section{Biseksi}

Jika $f$ kontinu pada $[a,b]$ dan

\begin{equation}
f(a)f(b)<0,
\end{equation}

maka sedikitnya satu akar berada di dalam interval.

Biseksi mengambil titik tengah

\begin{equation}
x_r=\frac{a+b}{2}.
\end{equation}

Kemudian interval diperbarui:

\begin{equation}
\begin{cases}
b\leftarrow x_r, & f(a)f(x_r)<0,\\
a\leftarrow x_r, & f(a)f(x_r)>0.
\end{cases}
\end{equation}

Galat iterasi yang dipakai:

\begin{equation}
e_a=
\frac{|x_{r,\mathrm{new}}-x_{r,\mathrm{old}}|}
{\max(1,|x_{r,\mathrm{new}}|)}.
\end{equation}

\subsection*{Langkah algoritma}

\begin{enumerate}
\item Pilih $a,b$ sehingga $f(a)f(b)<0$.
\item Hitung $x_r=(a+b)/2$.
\item Hitung $f(x_r)$ dan $e_a$.
\item Pilih interval baru berdasarkan tanda $f(a)f(x_r)$.
\item Ulangi sampai toleransi atau iterasi maksimum.
\end{enumerate}

\section{Regula Falsi}

Regula Falsi tetap memakai interval yang mengurung akar, tetapi kandidat akar tidak lagi menggunakan titik tengah.
Bangun garis yang melalui $(a,f(a))$ dan $(b,f(b))$.
Titik potong garis tersebut dengan sumbu-$x$ adalah

\begin{equation}
x_r
=
b-\frac{f(b)(b-a)}{f(b)-f(a)}
=
\frac{a f(b)-b f(a)}{f(b)-f(a)}.
\end{equation}

Setelah $x_r$ diperoleh, aturan pemilihan interval \textbf{sama dengan Biseksi}:

\begin{equation}
\begin{cases}
b\leftarrow x_r, & f(a)f(x_r)<0,\\
a\leftarrow x_r, & f(a)f(x_r)>0.
\end{cases}
\end{equation}

Jadi perbedaan inti kedua metode hanya terletak pada cara menghitung $x_r$.

\section{Titik Tetap}

Persamaan $f(x)=0$ diubah menjadi

\begin{equation}
x=g(x),
\end{equation}

kemudian dilakukan iterasi

\begin{equation}
x_{k+1}=g(x_k).
\end{equation}

Contoh:

\begin{equation}
e^{-x}-x=0
\qquad\Longrightarrow\qquad
x=e^{-x}.
\end{equation}

Syarat lokal sederhana untuk kecenderungan konvergen adalah

\begin{equation}
|g'(x^\ast)|<1.
\end{equation}

\section{Newton-Raphson}

Newton-Raphson menggunakan garis singgung pada titik $(x_k,f(x_k))$:

\begin{equation}
x_{k+1}
=
x_k-\frac{f(x_k)}{f'(x_k)}.
\end{equation}

Turunan harus ditentukan terlebih dahulu.

Contoh, untuk

\begin{equation}
f(x)=e^{-x}-x,
\end{equation}

diperoleh

\begin{equation}
f'(x)=-e^{-x}-1.
\end{equation}

\section{Secant}

Secant tidak memerlukan turunan analitik:

\begin{equation}
x_{k+1}
=
x_k-
f(x_k)
\frac{x_k-x_{k-1}}
{f(x_k)-f(x_{k-1})}.
\end{equation}

\section{Contoh kelas: Biseksi dan Regula Falsi}

Gunakan fungsi yang sama:

\begin{equation}
f(x)=x^3-x-1,
\qquad [a,b]=[1,1.5].
\end{equation}

Karena $f(1)=-1$ dan $f(1.5)=0.875$, interval mengurung akar.

Untuk Biseksi,

\begin{equation}
x_r=\frac{1+1.5}{2}=1.25.
\end{equation}

Untuk Regula Falsi,

\begin{equation}
x_r
=
1.5-\frac{0.875(1.5-1)}{0.875-(-1)}
\approx1.266667.
\end{equation}

Setelah nilai $x_r$ diperoleh, bagian kode yang memilih interval baru adalah identik.

\section{Program MATLAB yang dibahas di kelas}

Kode kelas hanya membandingkan Biseksi dan Regula Falsi.

\texttt{akar\_demo.m}

\lstinputlisting[language=Matlab]{matlab/akar_demo.m}

\section{Latihan / tugas utama}

\begin{enumerate}
\item \textbf{Biseksi.} Selesaikan $x^3-4x-9=0$. Tentukan interval sendiri, tampilkan tabel iterasi, simpan \texttt{errs}, dan gambar grafik konvergensi.
\item \textbf{Newton-Raphson.} Untuk $f(x)=e^{-x}-x$, tentukan $f'(x)$, gunakan sedikitnya dua tebakan awal, simpan \texttt{errs}, dan gambar grafik konvergensi.
\item \textbf{Secant.} Gunakan fungsi yang sama dengan Newton-Raphson, pilih dua tebakan awal, simpan \texttt{errs}, lalu bandingkan jumlah iterasi dan pola grafiknya dengan Newton-Raphson.
\end{enumerate}

\section{Latihan terintegrasi dari Lembar Kerja}

\begin{enumerate}
\item[\textbf{B1.1 (L1)}] Untuk \(f(x)=x^3-x-1\), buat tabel nilai \(f(x)\) untuk \(x=0.5,1,1.25,1.5,2\) dan tentukan interval yang masuk akal sebagai \textit{bracket}.
\item[\textbf{B1.2 (L2)}] Gunakan Biseksi untuk \(f(x)=x^3-x-1\) pada \([1,1.5]\). Lakukan 4 iterasi dan tuliskan tabel iterasinya.
\item[\textbf{B1.3 (L3)}] Gunakan Regula Falsi untuk fungsi yang sama dengan \(a=1\), \(b=1.5\). Lakukan 5 iterasi.
\item[\textbf{B1.4 (L4)}] Gunakan Newton-Raphson untuk \(f(x)=x^3-x-1\). Tentukan turunannya terlebih dahulu.
\item[\textbf{B1.5 (L5)}] Gunakan Newton-Raphson untuk \(f(x)=e^{-x}-x\). Tentukan turunannya, lakukan 5 iterasi, dan bandingkan konvergensinya dengan B1.4.
\item[\textbf{B2.1 (L1)}] Untuk \(e^{-x}-x=0\), gunakan \(g(x)=e^{-x}\) dan hitung 5 iterasi titik tetap dari \(x_0=0\).
\item[\textbf{B2.2 (L2)}] Evaluasi \(|g'(x^\ast)|\) di sekitar \(x^\ast\approx0.567\).
\item[\textbf{B2.3 (L3)}] Gunakan Secant untuk \(f(x)=e^{-x}-x\) dengan \(x_0=0\), \(x_1=1\), dan lakukan 5 pembaruan.
\item[\textbf{B2.4 (L4)}] Untuk \(\cos x-x=0\), tuliskan dua kemungkinan bentuk titik tetap dan diskusikan kecenderungan konvergensinya.
\item[\textbf{B2.5 (L5)}] Bandingkan Biseksi, Regula Falsi, Titik Tetap, Newton-Raphson, dan Secant dari kebutuhan tebakan awal, turunan, kecepatan, dan risiko divergensi.
\end{enumerate}
